% GENERATED FILE. Edit canonical lessons, then run npm run generate:books.
% canonical-source-hash: fnv1a64:e1064e115338fb54
% canonical-lessons: ZH-C345-jushuo, ZH-C345-congci, ZH-C345-qiyuan, ZH-C345-wangyou, ZH-C345-youbian, ZH-R345-review-341-345, ZH-R345-review-final-345

\chapter{Looking Back}
\label{ch:zh-looking-back}
\hlchaptermodality{345}

\begin{chapteropening}
\textbf{By the end of this chapter:} \emph{I can say what people say about the past, say what happened from then on, ask about origins and keep in touch.}

Complete the last lesson of chapter 345: I can say what people say about the past, say what happened from then on, ask about origins and keep in touch.
\end{chapteropening}

\section[jùshuō]{jùshuō — it is said that; reportedly}

\label{lesson:ZH-C345-jushuo}



\begin{quote}
Before the new one: say the Mandarin for a photo album, then the Mandarin for forever and ever.
\end{quote}

\subsection*{You'll want to know: jùshuō}
\emph{jùshuō} — \textquotedblleft{}it is said that; reportedly\textquotedblright{}.

\begin{quote}
\emph{jùshuō zhège chuánshuō yǒu liǎngqiān nián le} — They say this legend is two thousand years old.
\end{quote}

Its characters come later; for now it is a word for the ear.

\begin{cousinweb}
\emph{jùshuō} is built from \emph{jù} 'to rely on' and \zh{说} 'to say', both inherited from Middle Chinese, \emph{kjoH} and \emph{sywet} (Guangyun, in Baxter's transcription), which glosses the first 'to lean on' and 'to cite'.
\end{cousinweb}

\begin{grammarlens}[title={What you've built}]
One more word for telling what happened.
\end{grammarlens}

\subsection*{Your turn}
\begin{itemize}
\raggedright
  \item \emph{Say it:} \emph{jùshuō}
  \item \emph{Say it:} \emph{jùshuō}, once more
  \item \emph{Say it:} \emph{jùshuō zhège chuánshuō yǒu liǎngqiān nián le} — They say this legend is two thousand years old.
  \item \emph{Recall:} say the Mandarin for a photo album, then the Mandarin for forever and ever, then say \emph{jùshuō} again
\end{itemize}

\begin{tcolorbox}[breakable,colback=teal!4,colframe=teal!35!black,title={Before you move on}]
What does jùshuō mean? (\textquotedblleft{}It is said that; reportedly\textquotedblright{}.) Say it once more.
\end{tcolorbox}
\section[cóngcǐ]{cóngcǐ — from then on}

\label{lesson:ZH-C345-congci}



\begin{quote}
Before the new one: say the Mandarin for forever and ever, then the Mandarin for reportedly.
\end{quote}

\subsection*{You'll want to know: cóngcǐ}
\emph{cóngcǐ} — \textquotedblleft{}from then on\textquotedblright{}.

\begin{quote}
\emph{cóngcǐ wángzǐ hé gōngzhǔ zhù zài yìqǐ} — From then on the prince and the princess lived together.
\end{quote}

Its characters come later; for now it is a word for the ear.

\begin{cousinweb}
\emph{cóngcǐ} is built from \zh{从} 'from' and \emph{cǐ} 'this', both inherited from Middle Chinese, \emph{dzjowng} and \emph{tshjeX} (Guangyun, in Baxter's transcription), which glosses the first 'to follow' and the second with a word for 'to stop'; the sense 'this' is not given there.
\end{cousinweb}

\begin{grammarlens}[title={What you've built}]
One more word for telling what happened.
\end{grammarlens}

\subsection*{Your turn}
\begin{itemize}
\raggedright
  \item \emph{Say it:} \emph{cóngcǐ}
  \item \emph{Say it:} \emph{cóngcǐ}, once more
  \item \emph{Say it:} \emph{cóngcǐ wángzǐ hé gōngzhǔ zhù zài yìqǐ} — From then on the prince and the princess lived together.
  \item \emph{Recall:} say the Mandarin for forever and ever, then the Mandarin for reportedly, then say \emph{cóngcǐ} again
\end{itemize}

\begin{tcolorbox}[breakable,colback=teal!4,colframe=teal!35!black,title={Before you move on}]
What does cóngcǐ mean? (\textquotedblleft{}From then on\textquotedblright{}.) Say it once more.
\end{tcolorbox}
\section[qǐyuán]{qǐyuán — origin}

\label{lesson:ZH-C345-qiyuan}



\begin{quote}
Before the new one: say the Mandarin for reportedly, then the Mandarin for from then on.
\end{quote}

\subsection*{You'll want to know: qǐyuán}
\emph{qǐyuán} — \textquotedblleft{}origin\textquotedblright{}.

\begin{quote}
\emph{nǐ zhīdào Chūnjié de qǐyuán ma} — Do you know how the lunar new year began?
\end{quote}

Its characters come later; for now it is a word for the ear.

\begin{cousinweb}
\emph{qǐyuán} is built from \zh{起} 'to rise' and \emph{yuán} 'a source', both inherited from Middle Chinese, \emph{khiX} and \emph{ngjwon} (Guangyun, in Baxter's transcription), which says of the second that the head of a river is called so.
\end{cousinweb}

\begin{grammarlens}[title={What you've built}]
One more word for telling what happened.
\end{grammarlens}

\subsection*{Your turn}
\begin{itemize}
\raggedright
  \item \emph{Say it:} \emph{qǐyuán}
  \item \emph{Say it:} \emph{qǐyuán}, once more
  \item \emph{Say it:} \emph{nǐ zhīdào Chūnjié de qǐyuán ma} — Do you know how the lunar new year began?
  \item \emph{Recall:} say the Mandarin for reportedly, then the Mandarin for from then on, then say \emph{qǐyuán} again
\end{itemize}

\begin{tcolorbox}[breakable,colback=teal!4,colframe=teal!35!black,title={Before you move on}]
What does qǐyuán mean? (\textquotedblleft{}Origin\textquotedblright{}.) Say it once more.
\end{tcolorbox}
\section[wǎngyǒu]{wǎngyǒu — online friend}

\label{lesson:ZH-C345-wangyou}



\begin{quote}
Before the new one: say the Mandarin for from then on, then the Mandarin for the origin of a custom.
\end{quote}

\subsection*{You'll want to know: wǎngyǒu}
\emph{wǎngyǒu} — \textquotedblleft{}online friend\textquotedblright{}.

\begin{quote}
\emph{wǒ hé wǒ de wǎngyǒu jiànmiàn le} — I met my online friend in person.
\end{quote}

Its characters come later; for now it is a word for the ear.

Two third tones meet in this word: wǎng is said with a rising tone before yǒu, and the spelling keeps both marks.

\begin{cousinweb}
\emph{wǎngyǒu} is built from \zh{网} 'a net' and \emph{yǒu} 'a friend', both inherited from Middle Chinese, \emph{mjangX} and \emph{hjuwX} (Guangyun, in Baxter's transcription), which glosses the first 'a net for fish or game'. The internet sense is modern.
\end{cousinweb}

\begin{grammarlens}[title={What you've built}]
One more word for telling what happened.
\end{grammarlens}

\subsection*{Your turn}
\begin{itemize}
\raggedright
  \item \emph{Say it:} \emph{wǎngyǒu}
  \item \emph{Say it:} \emph{wǎngyǒu}, once more
  \item \emph{Say it:} \emph{wǒ hé wǒ de wǎngyǒu jiànmiàn le} — I met my online friend in person.
  \item \emph{Recall:} say the Mandarin for from then on, then the Mandarin for the origin of a custom, then say \emph{wǎngyǒu} again
\end{itemize}

\begin{tcolorbox}[breakable,colback=teal!4,colframe=teal!35!black,title={Before you move on}]
What does wǎngyǒu mean? (\textquotedblleft{}Online friend\textquotedblright{}.) Say it once more.
\end{tcolorbox}
\section[yóubiān]{yóubiān — postcode}

\label{lesson:ZH-C345-youbian}



\begin{quote}
Before the new one: say the Mandarin for the origin of a custom, then the Mandarin for an online friend.
\end{quote}

\subsection*{You'll want to know: yóubiān}
\emph{yóubiān} — \textquotedblleft{}postcode\textquotedblright{}.

\begin{quote}
\emph{nǐ jiā de yóubiān shì duōshao} — What's your postcode?
\end{quote}

Its characters come later; for now it is a word for the ear.

\begin{cousinweb}
\emph{yóubiān} is built from \emph{yóu} 'the post' and \emph{biān} 'to arrange', a short form of a longer name. Both are inherited from Middle Chinese, \emph{hjuw} and \emph{pjien} (Guangyun, in Baxter's transcription), which calls the first a lodge on the border; the dictionary lists the second under \emph{pen} and \emph{pjien}, both glossed 'to put in order'.
\end{cousinweb}

\begin{grammarlens}[title={What you've built}]
One more word for telling what happened.
\end{grammarlens}

\subsection*{Your turn}
\begin{itemize}
\raggedright
  \item \emph{Say it:} \emph{yóubiān}
  \item \emph{Say it:} \emph{yóubiān}, once more
  \item \emph{Say it:} \emph{nǐ jiā de yóubiān shì duōshao} — What's your postcode?
  \item \emph{Recall:} say the Mandarin for the origin of a custom, then the Mandarin for an online friend, then say \emph{yóubiān} again
\end{itemize}

\begin{tcolorbox}[breakable,colback=teal!4,colframe=teal!35!black,title={Before you move on}]
What does yóubiān mean? (\textquotedblleft{}Postcode\textquotedblright{}.) Say it once more.
\end{tcolorbox}
\section[(dialogue)]{Looking back: from Mooncakes to Looking Back}

\label{lesson:ZH-R345-review-341-345}



\begin{quote}
Say the Mandarin for sweet bean paste and a mooncake.
\end{quote}

\subsection*{Your turn}
\begin{itemize}
\raggedright
  \item \emph{Say it:} sweet bean paste, a mooncake, the filling of a dumpling, to taste food with care, trade in a shop
  \item \emph{Say it:} to enjoy the moon, the Mid-Autumn Festival, the time during an event, sweet osmanthus flowers, one's native place
  \item \emph{Say it:} bright moonlight, for many years now, deep in one's heart, good wishes sent to someone, to share good things
  \item \emph{Say it:} to keep something safe, a memory of the past, past events, a photo album, forever and ever
  \item \emph{Say it:} reportedly, from then on, the origin of a custom, an online friend, a postcode
  \item \emph{Read it:} \zh{饼} \zh{秋} — each one aloud, then a word that holds it
\end{itemize}

\begin{tcolorbox}[breakable,colback=teal!4,colframe=teal!35!black,title={Before you move on}]
Sweet bean paste? (\emph{dòushā}.) A mooncake? (\textbf{\zh{月饼}}, \emph{yuèbǐng}.) The filling of a dumpling? (\emph{xiànr}.) To taste food with care? (\emph{pǐncháng}.) Trade in a shop? (\textbf{\zh{生意}}, \emph{shēngyi}.) To enjoy the moon? (\emph{shǎngyuè}.) The Mid-Autumn Festival? (\textbf{\zh{中秋}}, \emph{Zhōngqiū}.) The time during an event? (\textbf{\zh{期间}}, \emph{qījiān}.) Sweet osmanthus flowers? (\emph{guìhuā}.) One's native place? (\emph{jiāxiāng}.) Bright moonlight? (\emph{yuèguāng}.) For many years now? (\textbf{\zh{多年}}, \emph{duōnián}.) Deep in one's heart? (\textbf{\zh{心中}}, \emph{xīnzhōng}.) Good wishes sent to someone? (\emph{zhùfú}.) To share good things? (\emph{fēnxiǎng}.) To keep something safe? (\emph{bǎocún}.) A memory of the past? (\emph{jìyì}.) Past events? (\emph{wǎngshì}.) A photo album? (\emph{xiàngcè}.) Forever and ever? (\emph{yǒngyuǎn}.) Reportedly? (\emph{jùshuō}.) From then on? (\emph{cóngcǐ}.) The origin of a custom? (\emph{qǐyuán}.) An online friend? (\emph{wǎngyǒu}.) A postcode? (\emph{yóubiān}.)
\end{tcolorbox}
\section[(dialogue)]{Looking Back — once more}

\label{lesson:ZH-R345-review-final-345}



\begin{quote}
Say the Mandarin for reportedly and from then on.
\end{quote}

\subsection*{Your turn}
\begin{itemize}
\raggedright
  \item \emph{Say it:} reportedly, from then on, the origin of a custom, an online friend, a postcode
\end{itemize}

\begin{tcolorbox}[breakable,colback=teal!4,colframe=teal!35!black,title={Before you move on}]
Reportedly? (\emph{jùshuō}.) From then on? (\emph{cóngcǐ}.) The origin of a custom? (\emph{qǐyuán}.) An online friend? (\emph{wǎngyǒu}.) A postcode? (\emph{yóubiān}.)
\end{tcolorbox}
